\subsection{Anforderungen an Wiederverwendbarkeit und Standardisierung}
\label{subsec:wiederverwendbarkeit-standardisierung}

Die Baseline soll gemeinsame Strukturen und öffentliche Arbeitsabläufe
vorgeben, aber kontrollierte Profilvarianten zulassen. Eine zentrale Pflege
vermeidet auseinanderlaufende Kopien. Tabelle~\ref{tab:anforderungen-template}
fasst die überprüfbaren Anforderungen zusammen.

\begin{table}[H]
  \centering
  \scriptsize
  \renewcommand{\arraystretch}{1.08}
  \caption{Zentrale Anforderungen an das Technologie-Template}
  \label{tab:anforderungen-template}
  \begin{tabularx}{\textwidth}{@{}p{0.08\textwidth}p{0.22\textwidth}p{0.39\textwidth}>{\raggedright\arraybackslash}X@{}}
    \toprule
    \textbf{ID} & \textbf{Anforderung} & \textbf{Operationalisiertes Ziel} & \textbf{spätere Prüfbasis} \\
    \midrule
    A-01 & Baseline & Grundstruktur muss generierbar sein. & Projektgenerierung \\
    A-02 & Struktur & Verantwortlichkeiten liegen an konsistenten Stellen. & Repository-Struktur \\
    A-03 & Profile & Nur vorgesehene Komponenten werden aktiviert. & Profilmodell \\
    A-04 & Workflows & Prüfen, Initialisieren, Installieren, Starten, Testen und Bauen besitzen gemeinsame Einstiegspunkte. & Tooling \\
    A-05 & Wartbarkeit & Gemeinsame Bestandteile werden zentral gepflegt. & Repository-Strategie \\
    A-06 & Nachvollziehbarkeit & Änderungen und Konfiguration sind versioniert und dokumentiert. & Versionsstand, Dokumentation \\
    A-07 & Erweiterbarkeit & Optionale Fähigkeiten besitzen dokumentierte Abhängigkeiten. & Capability-Modell \\
    \bottomrule
  \end{tabularx}
  \smallskip
  {\footnotesize Quelle: Eigene Darstellung.\par}
\end{table}

Die genannte Prüfbasis dient später der Bewertung dieser Soll-Kriterien.
\beginnerinput{workspace/statement/200_anforderungen/220}
